\section{讨论与结论}

\subsection{各改进对识别性能的贡献分析}

\subsubsection{动态特征的贡献}

从 13 维静态 MFCC 扩展到 39 维动态特征（MFCC + $\Delta$ + $\Delta^2$）
是 Advanced GMM 方法性能提升的关键因素之一。
静态 MFCC 仅反映当前帧的瞬时频谱包络，
而 $\Delta$ 和 $\Delta^2$ 分别编码了频谱随时间变化的速度和加速度。
这些动态特征在语音识别中尤为重要，
因为不同数字的发音在共振峰过渡、能量上升/衰减速率等方面存在显著差异。
例如，数字``6''和``9''在静态 MFCC 空间中可能较为接近，
但它们的频谱变化轨迹不同——``6''的发音通常以低频能量逐步衰减为特征，
而``9''则涉及更复杂的共振峰运动——这些差异在动态特征空间中更容易被区分。

\subsubsection{CMVN 的贡献}

CMVN（倒谱均值方差归一化）在句子级别将每维特征归零均值、单位方差，
有效抑制了信道卷积噪声和说话人差异带来的线性频谱畸变。
与 baseline 中使用的全局 StandardScaler 不同，
CMVN 是逐句独立进行的，因此能够自适应地补偿不同句子的录制条件差异，
而不依赖训练集的整体统计量。
这一处理使分类器更关注语音内容的相对变化模式，而非绝对频谱值。

\subsubsection{GMM 概率建模的优势}

GMM 分类器相较于 KNN 的核心优势在于其概率生成建模的特性：
\begin{enumerate}[label=(\arabic*)]
  \item \textbf{时序信息保留}：GMM 直接对变长帧序列中的所有帧进行联合概率建模，
        而非像 KNN 那样将所有帧压缩为两个统计量（均值和方差）。
        这使得模型能够利用语音中丰富的帧级频谱分布信息；
  \item \textbf{类内分布建模}：每个 GMM 通过 4 个高斯分量近似该类语音在特征空间中的多模态分布。
        考虑到同一数字在不同上下文中、由不同说话人发出时，
        其声学特征可能呈现多种模式（如清辅音段、元音段、过渡段），
        多分量 GMM 能够自然地捕捉这种多模态性；
  \item \textbf{软判决与概率校准}：GMM 基于对数似然进行决策，
        其输出可以被解释为样本属于各类的相对置信度，
        比 KNN 的硬最近邻规则具有更好的泛化能力；
  \item \textbf{小样本鲁棒性}：在小样本（每类仅 5 个训练样本）条件下，
        GMM 通过对所有帧进行密度估计来利用全部可用数据，
        而 KNN 仅依赖少数几个训练样本的距离排序，更容易受到个别异常帧的影响。
\end{enumerate}

\subsection{噪声鲁棒性分析}

噪声实验结果表明，Advanced GMM 方法的性能随 SNR 下降呈现明显的退化趋势，
但退化程度相较 baseline 有明显改善。
以 \qty{20}{dB} SNR 为例，GMM 平均准确率为 0.9125，而 baseline 仅为 0.5025。

从现象层面分析，噪声对 GMM 的影响主要体现在以下方面：
\begin{enumerate}[label=(\arabic*)]
  \item 加性噪声改变了每一帧的功率谱形状，
        导致 MFCC 及其差分特征发生系统性偏移，
        使测试帧在 GMM 似然空间中的分布偏离训练分布；
  \item 在低 SNR（如 \qty{-10}{dB}）下，
        语音信号几乎被噪声淹没，提取的特征主要反映噪声特性而非语音内容，
        导致所有类别的对数似然趋近，分类退化为近似随机猜测（准确率 $\approx$ 0.10，与 10 类随机概率一致）；
  \item 噪声对 $\Delta$ 和 $\Delta^2$ 特征的影响尤为显著——高频噪声使相邻帧间出现大幅随机波动，
        导致差分特征的信噪比低于静态 MFCC，在低 SNR 下反可能成为干扰。
\end{enumerate}

值得注意的是，在 \qty{5}{dB} 至 \qty{20}{dB} 区间，
GMM 准确率随 SNR 单调提升且标准差较小（$< 0.075$），
表明该方法在该 SNR 范围内具有相对稳定的噪声鲁棒性。

\subsection{与 Baseline 方法的性能对比}

表~\ref{tab:comparison_baseline_gmm} 给出了 Baseline KNN 与 Advanced GMM 在纯净和不同噪声条件下的性能对比。

\begin{table}[H]
  \centering
  \caption{Baseline KNN 与 Advanced GMM 性能对比}
  \label{tab:comparison_baseline_gmm}
  \begin{tabular}{@{}cccc@{}}
    \toprule
    SNR & Baseline KNN 准确率 & Advanced GMM 准确率 & 提升幅度 \\
    \midrule
    Clean & 0.8000 & \textbf{1.0000} & $+0.2000$ \\
    \qty{20}{dB} & 0.5025 & \textbf{0.9125} & $+0.4100$ \\
    \qty{10}{dB} & 0.2950 & \textbf{0.7250} & $+0.4300$ \\
    \qty{0}{dB}  & 0.3500 & \textbf{0.3800} & $+0.0300$ \\
    \bottomrule
  \end{tabular}
\end{table}

可以看出，Advanced GMM 在所有条件下均优于 baseline，
尤其在中高 SNR（\qty{10}{dB}--\qty{20}{dB}）区间提升幅度最大（约 $+0.41$--$+0.43$）。
在 \qty{0}{dB} 及以下极端噪声条件下，两种方法的差距缩小，
说明当噪声强度超过语音信号本身时，特征提取本身成为瓶颈，
分类器的改进空间有限。

\subsection{局限性与改进方向}

尽管 Advanced GMM 方法在纯净和小噪声条件下表现出色，但仍存在以下局限：

\begin{enumerate}[label=(\arabic*)]
  \item \textbf{对角协方差假设}：为简化参数估计，
        本方法使用对角协方差矩阵（\texttt{diag}），
        假设 39 维特征各维之间相互独立。
        这一假设在实际中并不完全成立——
        MFCC 相邻维度之间以及 MFCC 与 $\Delta$/$\Delta^2$ 之间存在相关性。
        使用全协方差或 tied 协方差可能进一步提升性能；
  \item \textbf{高斯分量数固定}：$K = 4$ 是根据经验选取的，
        未对不同分量数进行系统调优。
        分量数过少可能不足以捕获类内多模态分布，
        过多则在小样本下容易过拟合；
  \item \textbf{小样本训练集}：每类仅 5 个训练样本，
        GMM 的 4 个分量共需估计 $4 \times (39 + 39) + 3 \approx 315$ 个参数（对角协方差下，含均值、方差与混合权重），
        训练数据量相对有限，可能导致参数估计不够精确；
  \item \textbf{缺少有效语音帧筛选}：当前方法将所有帧（包括低能量静音帧）纳入 GMM 建模。
        在加噪条件下，静音段被填充为纯噪声帧，其分布与训练集中的静音帧分布不同，
        可能拉低整体对数似然；
  \item \textbf{噪声场景单一}：噪声鲁棒性测试仅使用一种混合噪声类型，
        未评估对不同噪声类型（如语音 babble、街道噪声、音乐噪声等）的泛化能力。
\end{enumerate}

未来改进方向可包括：
(1) 引入语音活动检测（VAD），仅对有效语音帧进行 GMM 建模和评分；
(2) 使用 Bayesian Information Criterion（BIC）或交叉验证自动选择最优高斯分量数；
(3) 尝试全协方差 GMM 或结合 UBM（Universal Background Model）的自适应方法；
(4) 在特征层面引入噪声抑制或谱减法预处理，提高低 SNR 下的特征质量。

% --- End of GMM discussion subsection ---
